\section{Современные вызовы и практическая реализация}
\label{sec:modern}

В XXI веке дипломная работа всё чаще выполняется в тесной связке с реальными индустриальными задачами. Это требует от студента не только теоретических знаний, но и владения современными инструментами разработки, навыков командной работы и умения документировать результаты.

\subsection{Интеграция науки и производства}
Программы ведущих университетов \cite{hse, mit} предусматривают выполнение ВКР на базе предприятий-партнёров. Такой подход позволяет решать актуальные задачи бизнеса и одновременно повышает качество подготовки выпускников. В таблице \ref{tab:partners} приведены примеры совместных проектов.

\begin{table}[htbp]
	\centering
	\caption{Примеры индустриальных партнёров и тематик дипломов}
	\label{tab:partners}
	\begin{tabular}{|p{3.5cm}|p{3.5cm}|p{5cm}|}
		\hline
		\textbf{Компания} & \textbf{Направление} & \textbf{Пример темы} \\
		\hline
		Яндекс & Машинное обучение & Разработка рекомендательной системы для образовательных курсов \\
		\hline
		ISPRAS & Системное программирование & Анализ уязвимостей в ядре Linux с помощью статического анализатора Svace \\
		\hline
		Сбер & Блокчейн & Прототип системы учёта дипломов на основе распределённого реестра \\
		\hline
	\end{tabular}
\end{table}

\subsection{Технологический стек современной ВКР}
Для успешной защиты выпускнику необходимо продемонстрировать владение современными средствами разработки. На листинге \ref{lst:c_example} показан пример программы на C, вычисляющей среднее арифметическое чисел из файла.

\begin{lstlisting}[caption={Фрагмент программы на C для анализа данных}, label=lst:c_example]
	#include <stdio.h>
	#include <stdlib.h>
	
	int main() {
		FILE *file = fopen("data.txt", "r");
		if (file == NULL) {
			printf("File open error\n");
			return 1;
		}
		
		int count = 0;
		double sum = 0.0;
		double value;
		
		while (fscanf(file, "%lf", &value) == 1) {
			sum += value;
			count++;
		}
		
		fclose(file);
		
		if (count > 0) {
			double mean = sum / count;
			printf("Number of elements: %d\n", count);
			printf("Arithmetic mean: %.2f\n", mean);
		} else {
			printf("File is empty or contains no numbers\n");
		}
		
		return 0;
	}
\end{lstlisting}

\subsection{Метрики качества дипломного проекта}
Современные требования включают не только техническую реализацию, но и документацию, тестирование, а также оценку экономической эффективности. На рисунке \ref{fig:metrics} представлена диаграмма значимости различных критериев на основе опроса работодателей \cite{career}.

\begin{figure}[htbp]
	\centering
	\includegraphics[width=0.7\textwidth]{example-image} % условное изображение
	\caption{Значимость критериев оценки дипломной работы (опрос 50 компаний)}
	\label{fig:metrics}
\end{figure}

Как показано в работах \cite{edutrends, hse}, наиболее важными факторами являются: применимость результата (35\%), качество кода (25\%), полнота документации (20\%) и инновационность (20\%).

\subsection{Заключение к главе}
Таким образом, дипломная работа в современной парадигме образования является не просто итоговым испытанием, а полноценным проектом, который готовит студента к реальной профессиональной деятельности. Применение передовых инструментов, сотрудничество с индустрией и ориентация на практический результат делают ВКР ключевым элементом образовательного процесса.